\documentclass[10pt,letterpaper]{article}
\usepackage{cornell-notes}
\title{Chapter 3: Field-Effect Transistors}
\author{}
\date{}
\setDocTitle{Chapter 3: Field-Effect Transistors}
\setDocSubtitle{Electronics}
\setDocOwner{Electronics Cornell Notes}
\setDocDate{\today}
\setCornellCollection{Electronics Cornell Notes}
\setCornellUnitType{Chapter}
\setCornellUnitNumber{3}
\setCornellUnitTitle{Field-Effect Transistors}
\hypersetup{pdftitle={Chapter 3: Field-Effect Transistors},pdfsubject={Cornell notes},pdfauthor={}}
\begin{document}
\maketitle
\begin{cornell}
\noterow{Core idea}{FET behavior follows gate control, channel conditions, and the operating region selected for switching or amplification.}
\noterow{Validation}{Check threshold spread, current, voltage, thermal behavior, and the intended switching or linear regime.}
\end{cornell}
\begin{summarybox}{Section summary}Choose a FET model that matches the device region and operating constraints.\end{summarybox}
\end{document}
